% Full helper-compiler proof of the 91-operation product-bound system.
\section{A positive product bound and two copied inputs: 91 operations}
\label{sec:product-bound}

The coefficient polynomial can be encoded as the difference of the two
already supplied codes. This removes the subtraction $q-C^2$ from the
product used in the third binary mask. Every negative coefficient then
creates an additional tested position. Two encoded copies of the input,
and positive complementary terms, make those extra tests exact without
adding arithmetic operations to the final system.

\begin{theorem}\label{thm:best91}
Every recursively enumerable subset of the positive integers has an
admissible fixed index $(V,H,T_L)$ for a system with 34 positive unknowns
and 22 equations admitting a straight-line certificate of 91 operations:
49 multiplications and 42 additions. The supplied inputs are
$(x,V,H,T_L)$, with three fixed components besides the positive queried
input $x$. Fixed numerals and equality tests are free.
\end{theorem}

\subsection{The changed source and its preliminary bounds}

Start from the system of Theorem~\ref{thm:best92}, retaining its
half-parameter Pell block and factored packed mask. Replace the coding
bound and product by
\begin{equation}\label{eq:product-bound-source}
 \ell+\sigma+\alpha=q,\qquad
 \sigma=(e-\ell)C^2,\qquad C=x+g.
\end{equation}
All supplied unknowns in these equations are positive integers. The
previous equations were
\[
 \ell+e+\alpha=q,\qquad
 \Omega=\lambda-e>0,\qquad
 \sigma=\Omega(q-C^2).
\]
The new system does not supply $\Omega$. Indeed, $C\ge2$ and
$\sigma>0$ imply $e-\ell>0$. If desired, a supplied positive
$\Omega=e-\ell$ and its free equality can be restored uniquely.
Their elimination therefore preserves the positive solution set.

Write $\Omega=e-\ell$ only as notation for this computed positive
integer. Equations~\eqref{eq:product-bound-source} give, before any
Pell or coefficient decoding,
\begin{equation}\label{eq:product-bound-ranges}
 \ell<q,\qquad \sigma<q,\qquad
 C^2\le\sigma<q,\qquad
 e=\ell+\Omega\le\ell+\sigma<q.
\end{equation}
The remaining definitions are
\begin{equation}\label{eq:product-bound-packing}
\begin{aligned}
 &H=H_0-3,\quad B=H+b+4=H_0+1+b,\quad
   \theta=B-4,\quad b=x+\beta,\\
 &\lambda(B-1)=q^2-1,\quad
   S_2=\ell+eq=V+t\theta,\quad n=q^8,\\
 &S=g+q^2(S_2+q^2\sigma),\\
 &T_+=q^2(1+\theta\lambda)+\ell(\theta q^4-b),\\
 &r=S(n^2-n)+T_+(n^2-1).
\end{aligned}
\end{equation}
Here $H_0$ is a sufficiently large fixed power of two, specified below.
In particular $S_2<q^2$ and $g<C<q$. Geometry gives $B\le q^2$;
the fixed threshold gives $3L\le B\le n$ and $\theta\lambda>b$.
The estimates used in Section~\ref{sec:affine-radix} consequently give
\begin{equation}\label{eq:product-bound-pre-pell}
 0<S,T_+<q^7<n,\qquad n^2-1\le r<2n^3.
\end{equation}
For example,
\[
 T_+>q^2\theta\lambda-b\ell>bq^2-bq>0,
 \qquad
 T_+<q^4(1+\theta\ell)<Bq^5\le q^7.
\]
The upper estimate uses $1+\theta\lambda<q^2$ and
$1+\theta\ell<Bq$. The bound on $S$ is stronger than its predecessor
because $\sigma<q$.

These are the same independent hypotheses used by the half-parameter
Pell proof in Section~\ref{sec:half-parameter-pell}. That proof now
applies in its stated order, before any coefficient test. It yields
the required central-binomial divisibility, $U=4^{2r+1}$, $q=B^L$,
and that $B,q,n$ are powers of two. Its sufficiency proof does not
assume even $r$. No discarded formula for $\sigma$ enters that Pell
argument.

\subsection{Nonnegative coordinates and two helpers}

Compile the represented finite Diophantine system into a circuit over
nonnegative integers, using the fresh-copy conventions of
Section~\ref{sec:affine-radix}. The actual positive input $x$ has
weight zero and is not split. Keep the homogeneous unit $\delta$ as
one physical coordinate. Every other logical value is a sum of three
distinct nonnegative physical coordinates. Different logical values
have disjoint physical groups.

For a nonnegative value $z$, the fixed forbidden-bit decomposition is
\begin{equation}\label{eq:product-bound-split}
 (z,0,0)\quad\hbox{if its $H_0$ bit is zero},\qquad
 (z-H_0,H_0/2,H_0/2)\quad\hbox{otherwise}.
\end{equation}
All pieces have that bit zero and are nonnegative. They can be fixed
before choosing an arbitrarily large power-of-two radix $B$.

Introduce two reserved logical helpers $V_0,V_1$, each the sum of its
own three physical coordinates. Every ordinary circuit row uses
$V_0$ in place of the input and mentions neither $x$ nor $V_1$.
Every monomial of such a row therefore has two non-input physical
factors. Fresh copies ensure that each nonzero physical square
coefficient is $\pm1$ and each nonzero physical cross coefficient is
$\pm2$. Dividing by the multiplicity in $C(T)^2$ gives coefficients
in $\{-1,1\}$.

The first four main rows are helper rows. For each $i=0,1$, use a
seed with base polynomial $-x^2$ and a reverse row with base
polynomial $x^2-V_i^2$. The seed will be augmented so that its actual
tested polynomial is $V_i^2-x^2$. The reverse row keeps its base
polynomial. This is the only intentional change to a main row caused
by the augmentation.

Include both signs of every ordinary zero row. In addition to the
circuit rows, include both signs of
\begin{equation}\label{eq:product-bound-normalization}
\begin{gathered}
 X^2-V_0^2,\qquad Y^2-V_0^2,\qquad Z^2-V_0^2,
 \qquad \delta'^2-\delta^2,\\
 2V_0X-2\delta\delta'-2u\delta.
\end{gathered}
\end{equation}
All these logical groups are distinct. Finish with three unpaired
$\delta^2$ rows. The last two receive the fivefold and sevenfold
padding described below. They have no negative coefficients.

\subsection{Weights and the positive compensation terms}

Let $m$ count the true non-input physical coordinates, including
$\delta$. Give them distinct weights
\begin{equation}\label{eq:product-bound-weights}
 v_i=8\cdot7^i\quad(0\le i<m),\qquad M=8\cdot7^{m-1}.
\end{equation}
After division by eight, a sum of at most six such weights has
base-seven digits below seven. Equality of any two such sums is
therefore equality of their multisets. In particular, a sum with
formal degree four cannot equal a sum with formal degree at most two.
The actual input has weight zero; its seed exception is treated
separately.

Let $s$ be the number of main rows, including the four helper rows
and the three final unit rows. Put
\begin{equation}\label{eq:product-bound-layout}
\begin{gathered}
 d_0=12M+8,\qquad
 t_j=(s+2)d_0+8M+jd_0\quad(0\le j<s),\\
 t_*=t_{s-1}=(2s+1)d_0+8M,\qquad K=t_*+3.
\end{gathered}
\end{equation}
For every divided base coefficient $a$ at a physical monomial of
weight $w$ in row $j$, put $aT^{t_j-w}$ in a complementary polynomial.

For each negative divided coefficient at $r=t_j-w$, choose a helper:
$V_1$ for an ordinary row, its own $V_i$ for seed $i$, and the other
$V_{1-i}$ for reverse row $i$. If its physical group is $P$, add
the six positive terms
\begin{equation}\label{eq:product-bound-augmentation}
 \sum_{h\le h',\ h,h'\in P}T^{r-v_h-v_{h'}}.
\end{equation}
The cross multiplicities in $C(T)^2$ recover the square of the sum
of the three helper coordinates. Mark every main target and every
negative base exponent as a tested start. Denote their set by
$\mathcal R$. In a seed, the negative exponent and main target are
the same start.

There are four relevant coefficient identities. First, an ordinary
or reverse augmentation has complementary weight
$w_{\rm negative}+w_{\rm helper}$ of degree four. It cannot affect
the main target, because a term of $C^2$ has degree at most two.
Second, at a negative position $r=t_j-w_{\rm negative}$, the negative
base term contributes $-x^2$. Another non-input base monomial could
contribute only if
\[
 w_{\rm base}=w_{\rm negative}+w_C.
\]
Degree and multiset uniqueness force $w_C=0$ and the identical base
monomial. The positive input square in a reverse row has weight zero
and cannot satisfy this equality.

Third, an augmentation contributes at that position only if
\begin{equation}\label{eq:product-bound-multiset}
 w_{{\rm negative},1}+w_{\rm helper}=w_{{\rm negative},2}+w_C.
\end{equation}
Every factor of a negative monomial in this row is disjoint from the
helper group. Both helper factors in~\eqref{eq:product-bound-multiset}
must consequently belong to the monomial on the right. Uniqueness
forces the identical negative pair and the identical helper pair.
Thus the extra tested polynomial is exactly
\begin{equation}\label{eq:product-bound-extra-target}
 V_{\rm helper}^2-x^2.
\end{equation}
Fourth, a seed has only $-x^2$ at its main target and its degree-two
helper complements. Its actual target is exactly $V_i^2-x^2$.

These arguments also exclude coefficient accumulation. Within a
non-seed augmentation, the partition into a negative pair and a
helper pair is unique by their disjoint groups. Seed helper pairs
have distinct weights. Base and augmented complementary exponents
have different degrees. Different row bands are separated by
$d_0>12M$. Thus all coefficients constructed so far remain $\pm1$.

\subsection{Resets, unit padding, and the two code polynomials}

Add a positive reset at every distinct tested start:
\begin{equation}\label{eq:product-bound-resets}
 D_{\rm reset}(T)=\sum_{r\in\mathcal R}T^{r-4}.
\end{equation}
Write $t_5=t_{s-2}$ and $t_7=t_{s-1}$ for the last two unit targets.
If $P(X)$ denotes the three physical indices of $X$, add
\begin{equation}\label{eq:product-bound-pads}
\begin{aligned}
 D_5(T)&=T^{t_5-1}
   +\sum_{h\in P(X)\cup P(Y)}T^{t_5-1-v_h},\\
 D_7(T)&=T^{t_7-1}
   +\sum_{h\in P(X)\cup P(Y)\cup P(Z)}T^{t_7-1-v_h}.
\end{aligned}
\end{equation}
Let $D(T)$ be the full complementary polynomial. Its base and
augmentation exponents are zero modulo eight, its resets are four
modulo eight, and its padding exponents are seven modulo eight.
These parts do not collide. Resets have distinct exponents, and the
padding terms are distinct within their disjoint final bands.
Consequently every nonzero coefficient of $D$ is still $\pm1$.

Define
\begin{equation}\label{eq:product-bound-codes}
 \ell_0(T)=\sum_{i=0}^{m-1}T^{v_i}
        +\sum_{r\in\mathcal R}(T^r+T^{r+1}+T^{r+2}),
 \qquad e_0(T)=\ell_0(T)+D(T).
\end{equation}
All indicator positions are distinct. Every negative coefficient of
$D$ occurs at a tested start and is canceled in $e_0$ by an indicator
digit one. Hence the digits of $\ell_0$ are zero or one and those of
$e_0$ are zero, one, or two. Both polynomials have zero constant
coefficient. The highest nonzero $D$ coefficient is the positive
padding term at $t_*-1$. Its leading $+1$ and remaining digits
$\pm1$ imply $D(B)>0$ for every integer $B\ge2$.

The indicator also permits dummy coordinate digits at the tested
positions. None can affect the low coefficient tests: the least
$D$ exponent is at least $t_0-4M$, the least dummy weight is at
least $t_0-2M$, and
\begin{equation}\label{eq:product-bound-dummies}
 2t_0-6M>t_*+2.
\end{equation}
Every true-coordinate weight lies below the least $D$ exponent.
Its raw coefficient and incoming carry are therefore zero.

Let $c_*$ count the input, all true physical coordinates, and all
permitted dummies, and put $D_1=\sum_h|[T^h]D(T)|$. Choose a power
of two $L>3K+2$ and then a power of two $H_0$ such that
\begin{equation}\label{eq:product-bound-fixed-threshold}
 H_0>\max\bigl(
   2\cdot4^{2L+1},\ 128\max(1,D_1)c_*^2,\ 3L,\ 1024\bigr).
\end{equation}
The supplied fixed index is
\begin{equation}\label{eq:product-bound-index}
 H=H_0-3,\qquad
 V=\ell_0(4)+4^Le_0(4),\qquad T_L=\psi_4(L).
\end{equation}
These choices precede the queried input and its witnesses.

\subsection{Canonical decoding after the Pell conclusions}

The bounds~\eqref{eq:product-bound-ranges} and the conclusion $q=B^L$
give $b\ell<Bq<q^2$. Thus
\begin{equation}\label{eq:product-bound-mask-blocks}
 T_+-1=(q^2-1-b\ell)+q^2\theta\lambda+q^4\theta\ell
\end{equation}
is the genuine block decomposition: its first two blocks are below
$q^2$, and the last is below $q^4$. The packed binomial condition
therefore gives the three independent binary masks for $g,S_2,\sigma$.

The middle mask is $\theta\lambda=(B-4)\sum_{j<2L}B^j$.
It requires every base-$B$ digit of $S_2$ to be at most three.
Write its digit polynomial as $F(T)$, with degree below $2L$.
Since $S_2\equiv V\pmod{B-4}$, one has
\[
 F(4)\equiv V\pmod{B-4}.
\]
Both sides are bounded by $4^{2L}-1$, while
\eqref{eq:product-bound-fixed-threshold} makes $B-4$ larger than
their possible difference. Hence $F(4)=V$. Uniqueness of the
base-four digit expansion gives
\[
 F(T)=\ell_0(T)+T^Le_0(T).
\]
Because $\ell,e<q$, the low and high $q$-blocks are their actual
values. Consequently $\ell=\ell_0(B)$ and $e=e_0(B)$.

The first mask now has digit $H_0=B-1-b$ at an indicator position
and digit $B-1$ elsewhere. It forces all non-indicator digits of
$g$ to vanish and clears the $H_0$ bit of each permitted digit.
The indicator has zero unit digit, so $g$ has zero unit digit.
Since $x<B$, adding $x$ introduces no carry, and $C=x+g$ is exactly
the physical coordinate polynomial evaluated at $B$.

Finally $\theta\ell_0(B)$ has digit $B-4$ at every indicator
position and zero elsewhere. The third mask says that each
corresponding digit of
\[
 \sigma=(e-\ell)C^2=D(B)C(B)^2
\]
is at most three. This includes every extra negative-coefficient
position and all three digits of each tested window.

\subsection{Signed carries and exact zero equations}

Every decoded coordinate digit is less than $B$, and $B\ge2H_0$.
Thus every raw coefficient of $D(T)C(T)^2$ has absolute value at most
\begin{equation}\label{eq:product-bound-cubic}
 D_1C(1)^2<D_1c_*^2B^2<B^3/256.
\end{equation}
The signed carry recurrence keeps every incoming carry in absolute
value below $B^2/128$. By~\eqref{eq:product-bound-dummies}, the only
raw residue classes through $t_*+2$ are zero, four, and seven
modulo eight.

At a tested start $r$, the positions $r-6,r-5$ are empty. Two
successive divisions by $B$ reduce the carry entering $r-4$ to
$-1$ or zero. Its raw reset polynomial is nonnegative and contains
$x^2$: the reset belonging to $r$ contributes $x^2$, and all other
reset contributions are positive. The other parts of $D$ have
different residue classes. Since $x\ge1$, the outgoing carry is
nonnegative and below $B^2$. The two empty positions $r-3,r-2$
therefore reduce it to zero.

At every start other than $t_5,t_7$, the position $r-1$ is also
empty. The following positions $r+1,r+2$ are always empty. Such a
test therefore sees its exact raw value modulo $B^3$. A negative
nonzero value satisfying~\eqref{eq:product-bound-cubic} has a top
digit greater than three and fails. The seed and reverse tests
thus give $V_i^2\ge x^2$ and $V_i^2\le x^2$, respectively. Since
these values are nonnegative,
\begin{equation}\label{eq:product-bound-helpers-fixed}
 V_0=V_1=x.
\end{equation}
All extra target values~\eqref{eq:product-bound-extra-target} now
vanish. Both signs of every ordinary row must pass, so every
ordinary zero equation is exact. In particular,
\begin{equation}\label{eq:product-bound-unit-guard}
 X=Y=Z=x,\qquad \delta'=\delta,\qquad
 x^2=\delta^2+u\delta,
 \qquad 0<\delta\le x<B.
\end{equation}
This establishes the input and guard constraints before using any
unit test.

At the two remaining preceding positions, the exact raw values
from~\eqref{eq:product-bound-pads} are
$x^2+2xX+2xY$ and $x^2+2xX+2xY+2xZ$. Their incoming carries
are zero. By~\eqref{eq:product-bound-unit-guard}, the three unit
tests are therefore the residues modulo $B^3$ of
\begin{equation}\label{eq:product-bound-three-units}
 \delta^2,\qquad
 \delta^2+\lfloor5x^2/B\rfloor,\qquad
 \delta^2+\lfloor7x^2/B\rfloor.
\end{equation}
There is no padding at an extra negative start: the two padded
rows contain no negative term, and other row bands are disjoint.

\subsection{The unit is one}

For completeness, the three-unit argument can be applied directly
to~\eqref{eq:product-bound-three-units}. Since $\delta<B$, the first
test gives
\[
 \delta^2=kB+s,\qquad k,s\in\{0,1,2,3\}.
\]
A square modulo four cannot have $s=2$ or $s=3$. If $s=1$, then
$\delta^2\le3B+3<B^2/16$, so $0<\delta<B/4$. The four square
roots of one modulo a power of two $B\ge8$ are represented by
$1,B/2-1,B/2+1,B-1$. Only the first lies in this range, giving
$\delta=1$.

Otherwise a nonunit solution would have $\delta^2=kB$ with
$1\le k\le3$ and $x^2\ge kB$. Write
\[
 h_c=\lfloor cx^2/B\rfloor=m_cB+r_c
 \quad(c=5,7),\qquad 0\le r_c<B.
\]
Because $x<B$, one has $h_c<cB$. The padded tests then force
$r_5,r_7\le3$ and $m_5,m_7\le3-k\le2$. On the other hand,
\[
 -7<7h_5-5h_7<5,\qquad |7r_5-5r_7|\le21.
\]
It follows that $|B(7m_5-5m_7)|<28$. Since $B>32$,
$7m_5=5m_7$, and their bounds imply $m_5=m_7=0$. This contradicts
$h_5\ge5k\ge5$ and $h_5=r_5\le3$. Thus $\delta=1$ in every
solution. The decoded circuit consequently represents the original
system at its actual input $x$, proving sufficiency.

\subsection{Necessity, including the auxiliary parity condition}

Given a represented solution, assign both helpers the value $x$,
set $\delta=\delta'=1$, $X=Y=Z=x$, and $u=x^2-1$, and supply all
circuit values and fresh copies. Split the non-input logical values
by~\eqref{eq:product-bound-split}; keep $\delta$ unsplit and set
all dummies to zero. Every ordinary, helper, and extra target has
raw value zero, and all three unit targets have raw value one.

After these finitely many physical values are fixed, choose $B$ to
be an arbitrarily large power of two. Require $B>H_0+1+x$, every
raw coefficient in absolute value below $B/4$, and $7x^2<B/4$.
Set
\begin{equation}\label{eq:product-bound-necessary-codes}
\begin{aligned}
 &b=B-H_0-1,\quad \beta=b-x,\quad g=C(B)-x,\\
 &\ell=\ell_0(B),\quad e=e_0(B),\quad q=B^L,\quad n=q^8,\\
 &\sigma=D(B)C(B)^2,\qquad \alpha=q-\ell-\sigma.
\end{aligned}
\end{equation}
The leading coefficient of $D$ makes $\sigma>0$. The encoded unit
makes $g>0$. Since $\deg(D C^2)<3K<L$, choosing $B$ still larger
makes $\alpha>0$. All forbidden-bit tests pass, and all middle-code
digits are at most two. Resets clear every signed carry. The two
padded unit carries are zero because $7x^2<B/4$. All third-mask
windows are therefore $(0,0,0)$ or $(1,0,0)$.

Geometry supplies a positive integer $\lambda$, and
\[
 t=\frac{\ell+eq-V}{B-4}>0
\]
is an integer: the numerator is the difference at $B$ and at four
of the nonconstant polynomial $\ell_0(T)+T^Le_0(T)$, which has
nonnegative coefficients. The three masks give the required binomial
divisibility for the newly determined $r$, whose growth bounds were
proved in~\eqref{eq:product-bound-pre-pell}.

Construct all non-auxiliary Pell witnesses as in
Section~\ref{sec:affine-radix}, with $U=4^{2r+1}$, the binomial
ratio value $Y_P$, the main and first Pell coordinates, and second
index $L$ with fixed $T_L=\psi_4(L)$. The half-parameter auxiliary
construction has one additional necessity hypothesis. Here $B$
divides both $g$ and $\ell$, since their unit digits vanish. Thus
both $S$ and $T_+$ in~\eqref{eq:product-bound-packing} are even,
and so is $r$. Consequently $J=2r+1$ is one modulo four.

Write $A=a+4$ and $D_P=A^2-1$, distinguishing the Pell coefficient
from the encoding polynomial $D(T)$. The positive construction in
Section~\ref{sec:half-parameter-pell} now applies with
\begin{equation}\label{eq:product-bound-auxiliary-witnesses}
\begin{gathered}
 m=2cJ,\qquad f=\chi_A(m),\qquad
 i=D_P\psi_A(m)/c^2,\qquad R=ic^2,\\
 y_{\rm aux}=\psi_R(J),\qquad
 u_{\rm aux}=\chi_R(J)/R,\\
 o=(u_{\rm aux}-c)/f,\qquad
 j=(u_{\rm aux}-J)/c.
\end{gathered}
\end{equation}
That proof establishes integrality and strict positivity using
$J\equiv1\pmod4$, and gives exactly the retained auxiliary norm
and congruence. All changed positive coding witnesses were supplied
in~\eqref{eq:product-bound-necessary-codes}. The difference $e-\ell$
is positive but is a computed register, not another supplied unknown.
This proves necessity for all 34 positive unknowns.

\subsection{Complete arithmetic certificate}

Relative to the 92-operation schedule, replace the three instructions
\[
 t_2=\lambda-e,\qquad b_\Sigma=\ell+e,\qquad
 S_3=t_2(q-C^2)
\]
by $t_2=e-\ell$, $b_\Sigma=\ell+\sigma$, and $S_3=t_2C^2$,
and delete the subtraction $q-C^2$. The supplied $\Omega$ and its
free equality are also deleted. Exactly one arithmetic operation
disappears, giving 49 multiplications and 42 additions.

The complete checker states all 22 source residuals afresh and
compares them with every certificate equality. If
$P_{\rm old}=(\lambda-e)(q-C^2)$ and $P_{\rm new}=(e-\ell)C^2$,
the aligned predecessor residuals change only by
\begin{equation}\label{eq:product-bound-source-differences}
 \sigma-e,\qquad
 -q^4(n^2-n)(P_{\rm new}-P_{\rm old}),\qquad
 P_{\rm new}-P_{\rm old},
\end{equation}
in the bound, packed-$r$, and product equations respectively.
The old positive-factor equation is omitted; the auxiliary congruence
is unchanged. To state the inherited residual corrections precisely,
let $F_j$ denote source residual $j$ in the checker's zero-based order,
and put $u_*=2r+1+jc$. The nonzero corrections are
\[
\begin{aligned}
 \mathcal E_2&=-\lambda F_3,\\
 \mathcal E_{16}&=F_{15}(u_*^2-y_{\rm aux}^2),\\
 \mathcal E_{17}&=F_3\bigl(\kappa+
       \rho(F_3-2(A-B))\bigr).
\end{aligned}
\]
Residuals $F_3$ and $F_{15}$ are independently exact geometry and
relaxed-norm equations. These corrections are therefore triangular
and vanish without circular reasoning. The generated schedule gives
all 91 primitive instructions and 22 free equality tests.

The general compiler and positive-domain proof received two
independent reviews. A supplementary sparse regression checks exact
symbolic target and padding identities in a sample with 35 true
coordinates, 25 main rows, 128 tested starts, 418 indicator digits,
and 976 complementary coefficients. Six canonical assignments and
four wide signed-carry assignments pass. These finite checks support
the general proof and do not assert exhaustive coverage of circuits
or integer witnesses. The construction proves an upper bound; it
does not claim that 91 is optimal.
